%% Question bank chapter file - generated by bank_tools.py build (do not edit while the bank is still being built:
%% it is regenerated from the per-test data in Question Bank\data). Every question is one bankq block:
%%   \begin{bankq}[id=..., source={...}, paper=.., qno=.., type=..., subtopic={...}, answer={...}, answer_src=key|solved,
%%                 working={...}, difficulty=1-10, diff_parts={K S I C F}, diff_note={...},
%%                 group=<passage id or empty>, also_in={duplicate copies in other tests}]
%% difficulty = min(10, K+S+I+C+F-1): K concepts 1-3, S steps 1-4, I insight 0-2, C computation 0-2, F format 0-1
%%   ...question LaTeX, options, figures (\bfig / \bfigi -> img/)...
%%   \end{bankq}
%% Passages shared by several questions are bankpassage blocks; their questions carry group=<passage id>.
\documentclass[10pt]{article}
\usepackage[a4paper,margin=1.8cm]{geometry}
\usepackage[T1]{fontenc}
\usepackage{amsmath,amssymb,mathtools}
\usepackage{graphicx}
\usepackage[export]{adjustbox}
\usepackage{xcolor,needspace,keyval}
\usepackage{enumitem}
\usepackage{booktabs,array,tabularx}
\usepackage[hidelinks]{hyperref}
\setlength{\parindent}{0pt}
\setlength{\parskip}{3pt}
\makeatletter
\define@key{bq}{id}{\def\bq@id{#1}}\define@key{bq}{source}{\def\bq@src{#1}}\define@key{bq}{paper}{\def\bq@paper{#1}}
\define@key{bq}{qno}{\def\bq@qno{#1}}\define@key{bq}{type}{\def\bq@type{#1}}\define@key{bq}{subtopic}{\def\bq@sub{#1}}
\define@key{bq}{answer}{\def\bq@ans{#1}}\define@key{bq}{answer_src}{\def\bq@asrc{#1}}\define@key{bq}{working}{\def\bq@work{#1}}
\define@key{bq}{group}{\def\bq@group{#1}}\define@key{bq}{also_in}{\def\bq@also{#1}}
\define@key{bq}{difficulty}{\def\bq@diff{#1}}\define@key{bq}{diff_parts}{\def\bq@dparts{#1}}\define@key{bq}{diff_note}{\def\bq@dnote{#1}}
\newenvironment{bankq}[1][]{\def\bq@work{}\def\bq@also{}\def\bq@group{}\def\bq@diff{?}\def\bq@dparts{}\def\bq@dnote{}\setkeys{bq}{#1}%
  \par\needspace{4\baselineskip}\medskip\noindent{\small\textbf{\bq@id}\hfill\color{gray}[\bq@src, Paper \bq@paper, Q\bq@qno]
  \bq@type\ $\cdot$ \bq@sub\ $\cdot$ \textbf{D\bq@diff}}\par\nopagebreak\relax}
  {\par{\small\color{teal}Answer: \textbf{\bq@ans}\ifx\bq@work\@empty\else\ (solved: \bq@work)\fi
  \ifx\bq@also\@empty\else\par\color{gray}Also in: \bq@also\fi}\par}
\define@key{bp}{id}{\def\bp@id{#1}}\define@key{bp}{source}{\def\bp@src{#1}}\define@key{bp}{questions}{\def\bp@qs{#1}}
\newenvironment{bankpassage}[1][]{\setkeys{bp}{#1}\par\needspace{4\baselineskip}\bigskip\noindent
  {\small\textbf{Passage \bp@id}\hfill\color{gray}[\bp@src; questions \bp@qs]}\par\nopagebreak\relax\itshape}{\par}
\makeatother
\newcommand{\bfig}[2]{\par\smallskip{\centering\includegraphics[width=#1pt,max width=0.62\linewidth,max height=0.32\textheight]{img/#2}\par}\smallskip}
\newcommand{\bfigi}[2]{\includegraphics[width=#1pt,max width=\linewidth,max height=0.32\textheight,valign=c]{img/#2}}
\newcommand{\bmissing}[1]{\textit{[figure missing in the source conversion: \detokenize{#1}]}}
\newenvironment{qoptions}{\begin{itemize}[leftmargin=1.8em,itemsep=1pt,topsep=2pt,label={}]}{\end{itemize}}
\newcommand{\opt}[2]{\item[(#1)] #2}
\providecommand{\fwstart}{}\providecommand{\fwend}{}\providecommand{\qfullsplit}{}
\newcolumntype{L}[1]{>{\raggedright\arraybackslash}p{#1}}
\newcolumntype{C}[1]{>{\centering\arraybackslash}p{#1}}
\newcolumntype{R}[1]{>{\raggedleft\arraybackslash}p{#1}}
\providecommand{\sectionhdr}[1]{}\providecommand{\subsectionhdr}[1]{}


\begin{document}
{\Large\bfseries Products of Vectors --- Question Bank}\par\smallskip
{\small Mathematics chapter M31. Questions from all converted test papers, grouped by source series. Answers: printed key where it matches the paper, otherwise solved (answer\_src=solved, with working). Exact repeats across papers appear once, with also\_in listing the other copies.}\par\medskip\hrule

\section*{Score-I papers}
\begin{bankq}[id=SI12-P2-MATH-S2-Q03, source={Score-I - Test 12 (09 Mar 2022)}, paper=2, qno=3, type=single, subtopic={Dot product}, answer={2000.00}, answer_src=key, difficulty=2, diff_parts={K1 S1 I1 C0 F0}, diff_note={concepts: Cauchy-Schwarz inequality | steps: S <= sqrt(9+16+100) * sqrt(16) = sqrt 2000 | I: Cauchy-Schwarz}]
Given $a^2+b^2+c^2\le16$ and $S=3a+4b+10c$ where $a,b,c\in R$. If the maximum value of $S$ is $\sqrt N$, then the value of $N$ is equal to
\end{bankq}

\section*{TOAS papers}
\begin{bankq}[id=TOAS17-P1-MATH-S1-Q04, source={TOAS - Test 17 (18 Sep 2022)}, paper=1, qno=4, type=multi, subtopic={Cross product}, answer={B,C}, answer_src=key, difficulty=4, diff_parts={K1 S2 I1 C0 F1}, diff_note={concepts: cross product and magnitude and angle | steps: b dot c; magnitudes; angle}]
$\vec a$ and $\vec b$ be two vectors such that $|\vec a|=1$, $|\vec b|=4$ and $\vec a\cdot\vec b=2$. If $\vec c=\left(2\vec a\times\vec b\right)-3\vec b$, then which of the following is/are correct ?
\begin{qoptions}
\opt{A}{$\vec b\cdot\vec c=48$}
\opt{B}{$\vec b\cdot\vec c=-48$}
\opt{C}{Angle between $\vec b$ and $\vec c$ is $\dfrac{5\pi}6$}
\opt{D}{Angle between $\vec b$ and $\vec c$ is $\dfrac\pi6$}
\end{qoptions}
\end{bankq}
\begin{bankq}[id=TOAS17-P1-MATH-S1-Q06, source={TOAS - Test 17 (18 Sep 2022)}, paper=1, qno=6, type=multi, subtopic={Scalar triple product}, answer={B,C,D}, answer_src=key, difficulty=4, diff_parts={K1 S2 I1 C0 F1}, diff_note={concepts: scalar triple product of linear combinations | steps: determinant of the coefficient matrix; multiplicative factor; volume}]
If volume of a parallelepiped determined by the vectors $\vec a,\vec b$ and $\vec c$ is 2, then volume of parallelepiped determined by the vectors
\begin{qoptions}
\opt{A}{$3(\vec a+\vec b)$, $(\vec b+\vec c)$ and $(\vec c+\vec a)$ is 24}
\opt{B}{$2(\vec a\times\vec b)$, $3(\vec b\times\vec c)$ and $(\vec c\times\vec a)$ is 24}
\opt{C}{$(\vec a-2\vec b)$, $(\vec b-2\vec c)$ and $(\vec c-2\vec a)$ is 14}
\opt{D}{$(3\vec a-\vec b)$, $(3\vec b-\vec c)$ and $(3\vec c-\vec a)$ is 52}
\end{qoptions}
\end{bankq}
\begin{bankq}[id=TOAS17-P1-MATH-S3-Q04, source={TOAS - Test 17 (18 Sep 2022)}, paper=1, qno=4, type=numerical, subtopic={Dot product}, answer={9.00}, answer_src=key, difficulty=3, diff_parts={K1 S1 I1 C0 F1}, diff_note={concepts: dot product with the circumcentre | steps: AO dot BC equals half the difference of squares of the sides}]
In a $\Delta ABC$, $AB=4$, $AC=5$, $O$ is circumcenter of $\Delta ABC$, then the value of $2\overrightarrow{AO}\cdot\overrightarrow{BC}$ ?
\end{bankq}
\begin{bankq}[id=TOAS17-P2-MATH-S1-Q02, source={TOAS - Test 17 (18 Sep 2022)}, paper=2, qno=2, type=single, subtopic={Dot product}, answer={D}, answer_src=key, difficulty=3, diff_parts={K2 S1 I1 C0 F0}, diff_note={concepts: vector identity from equal magnitudes, minimisation | steps: square both sides; dot product in terms of  | b}]
If $\vec a,\vec b$ are non-zero vectors such that $\left|\vec a+\vec b\right|=\left|\vec a-2\vec b\right|$, then -
\begin{qoptions}
\opt{A}{$\vec a\cdot\vec b=2\left|\vec b\right|^2$}
\opt{B}{$\vec a\cdot\vec b=\left|\vec b\right|^2$}
\opt{C}{Least value of $\vec a\cdot\vec b+\dfrac{4}{\left|\vec b\right|^2+2}$ is $2\sqrt2$}
\opt{D}{Least value of $\vec a\cdot\vec b+\dfrac{4}{\left|\vec b\right|^2+2}$ is $2\sqrt2-1$}
\end{qoptions}
\end{bankq}
\begin{bankq}[id=TOAS17-P2-MATH-S1-Q04, source={TOAS - Test 17 (18 Sep 2022)}, paper=2, qno=4, type=single, subtopic={Vector triple product}, answer={B}, answer_src=key, difficulty=5, diff_parts={K2 S2 I1 C1 F0}, diff_note={concepts: vector triple product, coplanar vectors | steps: expand the vector triple product; relate to a unit vector; angle condition; find c}]
Unit vectors $\vec a,\vec b$ and $\vec c$ are coplanar, unit vector $\vec d$ is perpendicular to them. If $\left(\vec a\times\vec b\right)\times\left(\vec c\times\vec d\right)=\dfrac{\hat i-2\hat j+2\hat k}{6}$ and angle between $\vec a$ and $\vec b$ is $30^\circ$ then $\vec c$ can be
\begin{qoptions}
\opt{A}{$\dfrac{\hat i-\hat j+\hat k}{\sqrt3}$}
\opt{B}{$\dfrac{\hat i-2\hat j+2\hat k}{3}$}
\opt{C}{$\dfrac{2\hat i-\hat j+\hat k}{\sqrt6}$}
\opt{D}{$\dfrac{\hat i-2\hat j+2\hat k}{6}$}
\end{qoptions}
\end{bankq}
\begin{bankq}[id=TOAS17-P2-MATH-S1-Q20, source={TOAS - Test 17 (18 Sep 2022)}, paper=2, qno=20, type=single, subtopic={Cross product}, answer={C}, answer_src=key, difficulty=4, diff_parts={K1 S2 I1 C1 F0}, diff_note={concepts: vector equation | steps: dot and cross products with the given vector; solve for v; magnitude}]
Let $\vec v$ be a vector which satisfies the condition $2\vec v+\vec v\times\left(\hat i+2\hat j\right)=2\hat i+\hat k$. If $3\left|\vec v\right|=\sqrt m$, then find the value of $m$ ?
\begin{qoptions}
\opt{A}{8}
\opt{B}{10}
\opt{C}{6}
\opt{D}{7}
\end{qoptions}
\end{bankq}
\begin{bankq}[id=TOAS17-P2-MATH-S2-Q05, source={TOAS - Test 17 (18 Sep 2022)}, paper=2, qno=5, type=numerical, subtopic={Cross product}, answer={15}, answer_src=solved, working={$(\vec r-\vec c)\times\vec b=0$ and $(\vec r-\vec c)\cdot\vec b=0$ give $\vec r=\vec c$, so $\vec a\cdot\vec r=8+7=15$}, difficulty=3, diff_parts={K1 S1 I1 C0 F1}, diff_note={concepts: vector equation with cross and dot products | steps: r - c parallel and perpendicular to b gives r = c; dot with a}]
Let $\vec a=2\hat i+\hat k$, $\vec b=\hat i+\hat j+\hat k$ and $\vec c=4\hat i+3\hat j+7\hat k$. If $\vec r\times\vec b=\vec c\times\vec b$ and $\vec r\cdot\vec b=\vec c\cdot\vec b$, then the value of $\vec a\cdot\vec r=$
\end{bankq}
\begin{bankq}[id=TOAS18-P1-MATH-S3-Q05, source={TOAS - Test 18 (30 Oct 2022)}, paper=1, qno=5, type=numerical, subtopic={Cross product}, answer={9}, answer_src=solved, working={Lines meet at $P(1,1,1)$; $\overrightarrow{AB}\times\overrightarrow{CD}=(1,-2,-2)$, $|{\cdot}|=3$, so $\overrightarrow{PQ}=\pm5(1,-2,-2)$ and $Q=(6,-9,-9)$, $\mu=9$}, difficulty=6, diff_parts={K2 S2 I1 C1 F1}, diff_note={concepts: intersection of lines, cross product | steps: intersection point; common perpendicular direction; position of Q; mu}]
The position vectors of two points $A$ and $C$ are $9\hat i-\hat j+7\hat k$ and $7\hat i-2\hat j+7\hat k$ respectively $\overrightarrow{AB}=4\hat i-\hat j+3\hat k$, $\overrightarrow{CD}=2\hat i-\hat j+2\hat k$, point of intersection of lines $AB$ and $CD$ is $P$. Vector $\overrightarrow{PQ}$ is perpendicular to $\overrightarrow{AB}$ and $\overrightarrow{CD}$ and $PQ=15$. If position vector of $Q$ is $6\hat i-9\hat j-\mu\hat k$ then value of $\mu$ is
\end{bankq}

\section*{TMAS papers}
\begin{bankq}[id=TMAS04-P1-MATH-S2-Q08, source={TMAS - Test 04 (19 Nov 2023)}, paper=1, qno=8, type=multi, subtopic={Vector triple product}, answer={A,C,D}, answer_src=key, difficulty=6, diff_parts={K1 S2 I2 C1 F1}, diff_note={concepts: vector triple product | steps: compute the cross product; solve for x y z; check}]
Let $\vec A=\hat i-2\hat j+3\hat k$, $\vec B=2\hat i+\hat j-\hat k$, $\vec C=\hat j+\hat k$. If the vector $\vec B\times\vec C$ can be expressed as $\vec B\times\vec C=x\vec A+y\vec B+z\vec C$ where $x,y,z$ are scalars, then which of the following is/are correct
\begin{qoptions}
\opt{A}{$x+2z=0$}
\opt{B}{$x=2y=2z=1$}
\opt{C}{$y+z=0$}
\opt{D}{$x=2y$}
\end{qoptions}
\end{bankq}
\begin{bankq}[id=TMAS04-P1-MATH-S2-Q11, source={TMAS - Test 04 (19 Nov 2023)}, paper=1, qno=11, type=multi, subtopic={Scalar triple product}, answer={A,D}, answer_src=key, difficulty=6, diff_parts={K1 S2 I2 C1 F1}, diff_note={concepts: coplanarity and vector products | steps: triple product conditions; solve; check}]
Which of the following option is/are correct ?
\begin{qoptions}
\opt{A}{If $\vec a=\alpha\hat i+2\hat j-3\hat k$, $\vec b=\hat i+2\alpha\hat j-2\hat k$ and $\vec c=2\hat i-\alpha\hat j+\hat k$ then $\left\{(\vec a\times\vec b)\times(\vec b\times\vec c)\right\}\times(\vec c\times\vec a)=0$ if $\alpha=\dfrac23$}
\opt{B}{If $\vec a=\alpha\hat i+2\hat j-3\hat k$, $\vec b=\hat i+2\alpha\hat j-2\hat k$ and $\vec c=2\hat i-\alpha\hat j+\hat k$ then $\left\{(\vec a\times\vec b)\times(\vec b\times\vec c)\right\}\times(\vec c\times\vec a)=0$ if $\alpha=\dfrac32$}
\opt{C}{If $\lambda\hat i+\hat j+2\hat k$, $\hat i+\lambda\hat j-\hat k$ and $2\hat i-\hat j+\lambda\hat k$ are coplanar then $\lambda$ is $-2$}
\opt{D}{If $\lambda\hat i+\hat j+2\hat k$, $\hat i+\lambda\hat j-\hat k$ and $2\hat i-\hat j+\lambda\hat k$ are coplanar then sum of all possible values of $\lambda$ is $0$}
\end{qoptions}
\end{bankq}
\begin{bankq}[id=TMAS04-P2-MATH-S3-Q01, source={TMAS - Test 04 (19 Nov 2023)}, paper=2, qno=1, type=numerical, subtopic={Cross product}, answer={6.00}, answer_src=key, difficulty=6, diff_parts={K1 S2 I2 C1 F1}, diff_note={concepts: area of a triangle from vector sides | steps: cross product of the two sides; magnitude; solve the equation}]
Let $\vec a$ and $\vec b$ are two non-zero, non collinear vectors ($|\vec a|=1$) such that vectors $3(\vec a\times\vec b)$ and $2(\vec b-(\vec a.\vec b)\vec a)$ represents two sides of a triangle. If area of triangle is $\dfrac34\left(|\vec b|^4+4\right)$, then find the value of $|\vec b|^2$
\end{bankq}
\end{document}
